\documentclass[11pt]{article}
\usepackage{rabbithole}

\phasenumber{7}
\phasetitle{Merging deeply}
\phasesubtitle{Strategies, why some merges make a commit, and teaching git to resolve a conflict once.}

\hypersetup{pdftitle={Rabbit Holes, Phase 7: Merging deeply},
            pdfauthor={Accelerate, Manipal University Jaipur},
            pdfsubject={git and GitHub handbook}}

\begin{document}
\maketitlepage

\section{What a merge actually does}

In the workshop you resolved a conflict by editing a file and running \co{git add}.
That is the whole user-facing experience of merging, and it hides a fair amount of
machinery worth understanding, because the machinery is what determines whether your
project's history is readable in a year.

A merge combines two branches. Git does it in one of two ways, and which one you get
depends entirely on the shape of the history.

\section{Fast-forward versus three-way}

\subsection{Fast-forward}

You branched off \co{main}, made two commits, and \co{main} has not moved since.

\begin{terminal}
A---B  main
     \
      C---D  feature
\end{terminal}

To merge \co{feature} into \co{main}, git does not need to combine anything. It just
moves the \co{main} label forward to \co{D}.

\begin{terminal}
A---B---C---D  main, feature
\end{terminal}

That is a \textbf{fast-forward}. No merge commit is created, because nothing was
merged. The history becomes a straight line, and afterwards there is no record that
a branch ever existed.

\subsection{Three-way merge}

Now \co{main} has moved on while you worked.

\begin{terminal}
A---B---E---F  main
     \
      C---D  feature
\end{terminal}

There is no way to move a label and get a correct answer. Git finds the \textbf{merge
base}, the most recent commit both branches share, which is \co{B}. Then it compares
three points: the base, and each of the two tips. Hence three-way.

Where only one side changed a region, git takes that change. Where both changed the
same region differently, it stops and asks you. That is a conflict, and it is the
only situation in which git ever asks.

The result is a \textbf{merge commit}, the only kind of commit with two parents:

\begin{terminal}
A---B---E---F---M  main
     \         /
      C---D---'
\end{terminal}

\begin{why}
Understanding that conflicts only happen when \emph{both} sides changed the same
region explains why the Phase 3 exercise worked. Everyone inserted a name at the
same position in the credits block, so every branch after the first had both sides
changing the same region. If people had added names in different places, git would
have merged them silently.

It also explains the Phase 2 filename rule. Different files never share a region, so
they can never conflict.
\end{why}

\subsection{Controlling which one you get}

\begin{shell}
git merge --no-ff feature
\end{shell}

Forces a merge commit even when a fast-forward was possible. The branch stays
visible in the history as a distinct piece of work.

\begin{shell}
git merge --ff-only feature
\end{shell}

Refuses to merge unless it can fast-forward. Useful in scripts and as a safety
check: it fails rather than silently creating a merge you did not intend.

\begin{shell}
git merge --squash feature
\end{shell}

Takes all the changes, stages them, and makes no commit. You commit once yourself.
The branch's individual commits are not recorded at all. This is what GitHub's
``Squash and merge'' button does.

\section{Which button on GitHub}

Every pull request offers three merge options, and teams argue about them. They map
directly onto the above.

\vspace{0.5em}
\begin{tabularx}{\linewidth}{@{}L{3.4cm} X@{}}
\toprule
\rhtablehead{Option} & \rhtablehead{What you get} \\
\midrule
Create a merge commit & Every commit from the branch, plus a merge commit. Full
history, more noise. \\
Squash and merge & One commit on \co{main} containing everything. Clean history,
individual steps lost. \\
Rebase and merge & Each commit replayed onto \co{main}, no merge commit. Linear
history, hashes change. \\
\bottomrule
\end{tabularx}
\vspace{0.5em}

\begin{note}
For a workshop, and for most student projects, \textbf{squash and merge} is the
right default. Contributors are still learning to make tidy commits, and squashing
means a messy branch still produces one clean entry on \co{main}.

For a project where contributors write careful, self-contained commits, squashing
throws away real information, and a merge commit or rebase preserves it.

There is no universally correct answer, which is why the argument never ends. Pick
one, write it in \co{CONTRIBUTING.md}, and stop discussing it.
\end{note}

\section{Merge strategies}

The strategy is the algorithm git uses to combine the trees.

\vspace{0.5em}
\begin{tabularx}{\linewidth}{@{}L{2.9cm} X@{}}
\toprule
\rhtablehead{Strategy} & \rhtablehead{Use} \\
\midrule
\co{ort} & The default since Git 2.34. Handles renames and directory moves far
better than what it replaced. You will almost never specify it. \\
\co{recursive} & The old default. Still available, no longer necessary. \\
\co{resolve} & Older and simpler, handles only two heads. \\
\co{octopus} & Merges more than two branches at once. Refuses if there are
conflicts. Rare. \\
\co{ours} & Discards the other branch's changes entirely, keeps the merge commit as
a record. \\
\co{subtree} & For merging a project into a subdirectory of another. \\
\bottomrule
\end{tabularx}
\vspace{0.5em}

\begin{shell}
git merge -s ours old-branch
\end{shell}

\subsection{Strategy options, which are different}

This is a genuine trap. \co{-s ours} is a strategy. \co{-X ours} is an option
\emph{to} a strategy, and it means something much narrower.

\begin{shell}
git merge -X ours feature
\end{shell}

This does a normal merge, taking every non-conflicting change from both sides, and
resolves only the \emph{conflicting} regions in favour of your side.

\begin{trap}
\co{-s ours} throws away everything from the other branch. \co{-X ours} keeps
everything that did not conflict.

The letter is the only difference and the outcomes are wildly different. If you mean
``take their version when we disagree'' you want \co{-X theirs}. If you type
\co{-s theirs} you will get an error, because it does not exist, which is the one
piece of luck in this design.
\end{trap}

\section{Better conflict markers}

The default conflict display shows you two versions. There is a better one.

\begin{shell}
git config --global merge.conflictStyle zdiff3
\end{shell}

Now conflicts include the original, from before either side changed it:

\begin{terminal}
<<<<<<< HEAD
timeout = 30
||||||| base
timeout = 10
=======
timeout = 60
>>>>>>> feature
\end{terminal}

The middle section changes everything. Without it you see two numbers and have to
guess the intent. With it you can see that one side tripled the timeout and the
other sextupled it, both were raising it from ten, and the answer is probably sixty
rather than a random pick.

\co{diff3} is the older version of the same idea. \co{zdiff3} is tidier and
available from Git 2.35 onwards. There is no reason not to turn it on.

\section{rerere}

\textbf{Re}use \textbf{re}corded \textbf{re}solution. Git watches how you resolve a
conflict and, if it ever sees the identical conflict again, resolves it the same way
without asking.

\begin{shell}
git config --global rerere.enabled true
\end{shell}

That is the entire setup. From then on it works quietly.

\begin{why}
This sounds like a small convenience until you have a long-lived branch that you
rebase onto a moving \co{main} every day. Without \co{rerere} you resolve the same
three conflicts every single time, identically, forever. With it you resolve them
once.

It also helps in the ordinary case of a merge that you abort, reconsider, and redo.
Git remembers what you decided the first time.
\end{why}

Inspect what it has stored:

\begin{shell}
git rerere status
git rerere diff
git rerere forget <path>
\end{shell}

\co{forget} is what you want when you resolved something wrongly and git is now
cheerfully repeating your mistake.

\section{Merge tools}

For a conflict spanning many regions, editing markers by hand gets tiring.

\begin{shell}
git mergetool
\end{shell}

This opens a three-pane view: their version, the merged result, yours. Configure
which program it uses:

\begin{shell}
git config --global merge.tool vscode
git config --global mergetool.vscode.cmd "code --wait $MERGED"
git config --global mergetool.keepBackup false
\end{shell}

That last line stops git leaving \co{.orig} files everywhere, which is otherwise a
persistent small annoyance.

\begin{note}
Editors have quietly made this obsolete for most people. VS Code, JetBrains IDEs and
Neovim all detect conflict markers and offer ``accept current, accept incoming,
accept both'' inline. That is usually faster than a separate three-pane tool.

Knowing what the markers mean still matters, because the buttons are just editing
text and you need to know when ``accept both'' is wrong.
\end{note}

\section{Branching workflows}

The strategies above are mechanics. A workflow is the social agreement about which
branches exist and what may be merged into them.

\subsection{GitHub flow}

One long-lived branch, \co{main}, which is always deployable. Every change is a short
branch that becomes a pull request and is merged and deleted.

That is the whole workflow. It is what this repository uses, what most web projects
use, and what you should default to.

\subsection{Trunk-based development}

Everyone commits to \co{main}, many times a day, with branches living hours rather
than days. Unfinished work is hidden behind feature flags rather than in branches.

Requires strong automated testing, because there is no branch to hold broken work.
Used heavily at large companies where long-lived branches cause more pain than
feature flags do.

\subsection{Git flow}

Five kinds of branch: \co{main}, \co{develop}, \co{feature/*}, \co{release/*},
\co{hotfix/*}. Designed in 2010 for software shipped as versioned downloads.

\begin{note}
Git flow is widely taught and, for most projects today, the wrong choice. Its author
has since added a note to the original article saying it was designed for a world
where you release versioned software on a schedule, and that continuously deployed
web applications should use something simpler.

If you inherit it, fine. Do not adopt it for a student project. The overhead is real
and the benefit only appears when you genuinely maintain several released versions
at once.
\end{note}

\subsection{Merge queues}

On a busy repository, a pull request can be tested against a \co{main} that changes
before it merges, so a green PR can still break \co{main}.

A merge queue solves this: pull requests enter a queue, each is tested against the
result of everything ahead of it, and merges only if that passes. GitHub has this
built in, under branch protection settings. It becomes worth the setup somewhere
around the point where merge collisions happen weekly.

\section{Reference}

\vspace{0.5em}
\begin{tabularx}{\linewidth}{@{}L{6.2cm} X@{}}
\toprule
\rhtablehead{Goal} & \rhtablehead{Command} \\
\midrule
Always record the branch & \co{git merge -{}-no-ff <branch>} \\
Refuse anything but a fast-forward & \co{git merge -{}-ff-only <branch>} \\
Combine into one commit yourself & \co{git merge -{}-squash <branch>} \\
Prefer my side on conflicts only & \co{git merge -X ours <branch>} \\
Prefer their side on conflicts only & \co{git merge -X theirs <branch>} \\
Discard their branch, keep the record & \co{git merge -s ours <branch>} \\
See the original in conflicts & \co{git config merge.conflictStyle zdiff3} \\
Never resolve the same conflict twice & \co{git config rerere.enabled true} \\
Find where two branches diverged & \co{git merge-base main feature} \\
Back out of a merge & \co{git merge -{}-abort} \\
\bottomrule
\end{tabularx}
\vspace{0.5em}

\checkyourself{
\begin{enumerate}
  \item When does git fast-forward instead of making a merge commit?
  \item What are the three inputs to a three-way merge?
  \item Explain the difference between \co{-s ours} and \co{-X ours}.
  \item What does \co{zdiff3} add to a conflict marker, and why is it useful?
  \item Your branch conflicts with \co{main} the same way every day. What do you
    turn on?
\end{enumerate}}

\vspace{1em}
{\color{rhborder}\rule{\linewidth}{0.8pt}}

\textbf{Next:} Phase 8, \emph{Remotes and big repositories}, covering what a refspec
is, how to work on several branches at once without stashing, and what to do when
the repository is too large to clone comfortably.

\end{document}
